\documentclass{report}
\usepackage{amsmath,amssymb}
\setlength{\parindent}{0mm}
\setlength{\parskip}{1em}
\begin{document}
\begin{center}
\rule{6in}{1pt} \
{\large Gunay Dogan \\
{\bf Fast iterative algorithm for shape reconstruction in image processing}}

Theiss Research \\ National Institute of Standards and Technology \\ 100 Bureau Drive \\ Stop 8910 \\ Gaithersburg \\ MD 20899-8910
\\
{\tt gunay.dogan@nist.gov}\end{center}

Many image processing tasks are naturally expressed as shape optimization
problems, in which the main goal is to reconstruct a set of shapes, such
as curves in 2d or surfaces in 3d, based on some given data. Examples of
such tasks are image segmentation, the task of identifying the objects or
uniform regions in a given image, or certain inverse problems, where the
target is not the material medium itself, but the structures or anomalies
in the medium. Often, at the outset of such problems, one has the data at
hand and some rough idea of what the shapes should or should not look
like, but does not have a direct way to compute the shapes from the data.
In such a situation, an effective approach is to define a shape energy
and to compute shapes that minimize this problem-specific shape energy.
The shape energy would typically have a data fidelity component minimized
by the correct shapes (based on the physics that relates the shape to the
data), also a smoothness term, such a length or area penalty that favors
smoother shapes over nonsmooth shapes.

In this work, we propose an efficient and robust iterative algorithm to
compute the solutions of such shape reconstruction problems in image
processing. At each iteration of the algorithm, we compute a gradient
descent velocity and use it to deform a candidate shape in a way that
keeps decreasing its energy until the algorithm stops at a minimum of the
energy. The starting point of our method is a continuous shape
optimization formulation and analytical computation of the first and
second shape derivatives based on the energy formulation. This
information enables us to devise Newton-type gradient descent methods
that converge in reduced number of iterations and significantly reduce
the actual computation time. Moreover we are able to impose the choice of
descent velocities to be in $H^1$ Sobolev spaces defined on the shape and
perform the gradient descent with smoother velocities. This proves to be
advantageous in the case of noisy and incomplete data; it reduces the
computation time and helps preserve a quality representation.

We discretize the continuous algorithm using the finite element method.
Unlike typical partial differential equation (PDE) applications, we need
to discretize both the geometry (which is also an unknown) and the
continuous equations (PDEs and integrals) associated with the gradient
descent. As the meshes representing the shapes and regions move and
deform through the iterations, we need to continually update and maintain
the mesh in order to preserve accuracy. Geometric accuracy, as well as
numerical accuracy, is achieved by maintaining a few errors indicators,
the total of which is used to ensure that we have a valid gradient
descent direction at each iteration. These error indicators, in
conjunction with geometric mesh adaptivity, are used to tune the accuracy
of the method at each iteration and are key to the economy of the
algorithm. In particular, we employ a coarse-to-fine continuation
approach and impose accuracy gradually as we get close to the solution,
thus save in computation when we are far away from the optimum. Putting
these ingredients together results in an effective algorithm applicable
to many shape reconstruction problem in imaging. We demonstrate the power
of the method with several examples.


\end{document}
